\chapter{Introducción}\label{cap:intro}

\section{El problema: un médico que no puede esperar a la Tierra}

En un viaje a Marte, una consulta con el cirujano de vuelo tarda entre 6 y 44 minutos de ida y vuelta, lo que tarda la
luz, y durante la conjunción solar el enlace se corta por semanas. Una emergencia médica no puede esperar eso. Por eso
el desafío \enquote{Create Health Monitoring Software for Astronauts on Space Missions} del NASA International Space
Apps Challenge 2026 pide monitoreo, detección y primera respuesta \emph{dentro} de la nave.

NOVA es la respuesta del equipo: un copiloto de salud que vive entero a bordo y funciona sin internet. Su núcleo de
vuelo usa F~Prime, el software de vuelo de NASA JPL; recibe los signos vitales que cada tripulante registra a diario en
una estación médica física y en su teléfono, aplica umbrales fijos y revisables y levanta alertas. Cinco modelos
locales (lenguaje, búsqueda, oído, voz y visión) explican lo que pasa y guían los protocolos de emergencia. La regla
del proyecto es estricta: lo crítico vive en reglas fijas y archivos revisables, y el modelo de lenguaje explica, no
decide.

Este informe trata de una pieza de NOVA: cómo saber qué es normal para cada persona cuando el cuerpo cambia en el
espacio.

\section{Por qué hace falta transferir lo aprendido en Tierra}

En la Tierra, la forma más simple de vigilar a una persona es compararla consigo misma: su media de los últimos días,
más o menos unas desviaciones. Antes del vuelo hay semanas de mediciones así. El problema es que la microgravedad
cambia el punto de partida. Los fluidos se desplazan hacia la cabeza, el corazón trabaja de otra forma, se pierde
masa muscular y ósea, y el peso baja \parencite{norsk2020,hughson2018,pavy2007}. Un monitor afinado en Tierra empieza
a ver anomalías donde solo hay adaptación, y una tripulación que recibe alarmas todos los días deja de mirarlas.

Lo que se necesita es una \enquote{nueva normalidad} personal: el valor que se espera de cada tripulante en cada día de
misión, con un rango que avise cuando algo se sale de lo esperable para \emph{esa} persona en \emph{ese} momento. Para
construirla hay muy pocos datos del propio tripulante en el espacio, sobre todo al principio. Es el problema clásico
del transfer learning: aprovechar lo que se aprendió en un dominio con muchos datos (la Tierra) para un dominio
parecido, pero distinto, en el que hay pocos (el espacio) \parencite{pan2010}.

\section{Preguntas de este trabajo}

El trabajo responde tres preguntas, que corresponden a los tres gráficos centrales del proyecto:

\begin{enumerate}[label=\textbf{P\arabic*.}, leftmargin=2.2em]
\item ¿Sirve en el espacio el monitor de Tierra? Si no sirve, ¿cuántas falsas alarmas da y en qué variables?
\item ¿Cuánto mejora un rango esperado adaptado, con las mismas falsas alarmas? ¿Cuántos días propios necesita y qué
  aporta lo aprendido del reposo de otras personas?
\item ¿El pre-entrenamiento en Tierra mejora el pronóstico a 7 días de un modelo de series de tiempo, frente a
  entrenarlo desde cero y frente a una línea de base simple?
\end{enumerate}

Las preguntas se responden con dos tareas. La primera, el \emph{rango esperado}, es un modelo residual pequeño: la
línea base de Tierra, un corrector aprendido de otros sujetos y los primeros días de la propia persona. La segunda, el
\emph{pronóstico}, usa PatchTST \parencite{nie2023}, un transformer para series de tiempo de la librería
\texttt{transformers} de Hugging Face \parencite{wolf2020}, pre-entrenado en Tierra y ajustado al espacio, con
Chronos-2 \parencite{ansari2025} como referencia sin adaptar.

\section{Alcance y advertencias}

\begin{cuidado}
Los datos del \enquote{espacio} son de reposo en cama con la cabeza inclinada a $-6^\circ$, el análogo terrestre más
usado de la microgravedad \parencite{pavy2007,hargens2016}. No son datos de vuelo. Las anomalías que se usan para medir
la detección son escalones sintéticos inyectados sobre los datos reales, y sus magnitudes son supuestos, salvo la del
pulso. La conexión con la nave está diseñada, pero no construida.
\end{cuidado}

El trabajo cubre los primeros 2 a 3 meses de una misión, que es lo que duran las campañas de reposo en cama (entre
\cifra{datos.chk.cama.min} y \cifra{datos.chk.cama.max} días en cama según el sujeto). No cubre radiación, CO$_2$ de
cabina ni ánimo: el reposo en cama no los tiene. La radiación queda diseñada como entrada futura del corrector.

\section{Cómo leer este informe}

El \cref{cap:historia} cuenta la historia del transfer learning y del ajuste fino desde sus orígenes en la psicología de
principios del siglo~\textsc{xx} hasta los modelos fundacionales de series de tiempo, y deja la teoría que se usa
después. El \cref{cap:datos} describe los datos y lo que cambia en cama. El \cref{cap:metodo} explica los dos modelos, la
validación y cómo se evitó hacerse trampa. El \cref{cap:resultados} presenta los resultados; el \cref{cap:auditoria}, la
auditoría independiente del tercer run; el \cref{cap:discusion}, qué significan a la luz de la literatura y cuáles son
sus límites; el \cref{cap:nova}, cómo se conecta con la nave y qué viene en la corrida siguiente, y el
\cref{cap:conclusiones}, las respuestas a las tres preguntas. Los apéndices traen las 26 decisiones que se tomaron
sobre el trabajo, la reproducibilidad y tablas de apoyo.

\begin{clave}
En los gráficos, el \textcolor{serieB}{\textbf{naranjo}} es siempre el monitor de Tierra, el
\textcolor{serieA}{\textbf{azul}} es el modelo adaptado y el \textcolor{serieC!80!black}{\textbf{verde}} es el monitor
de Tierra con el ancho calibrado o, en la curva de aprendizaje, el modelo que solo usa los días propios. En los
gráficos del entrenamiento las series llevan su nombre al lado. Las cifras de resultados salen de la salida del kernel
por un script, así que el texto y las tablas no pueden discrepar.
\end{clave}
